\documentclass[fleqn,usenatbib]{mnras}

\usepackage[T1]{fontenc}
\usepackage{newtxtext,newtxmath}
\usepackage{graphicx}
\usepackage{amsmath}
\usepackage{amssymb}
\usepackage{booktabs}
\usepackage{longtable}
\usepackage{multirow}
\usepackage{xcolor}

\title[TESS Exoplanet Pipeline]{TESS Exoplanet Pipeline: Sphinx Documentation and an MNRAS-Style Scientific Analysis Framework}

\author[Siya]{Siya\\
Department of Astronomy and Astrophysics\\
\texttt{siya@example.com}}

\begin{document}

\maketitle

\begin{abstract}
We present the initial manuscript scaffold for the TESS exoplanet pipeline project. The paper will describe the pipeline architecture, detection workflow, stellar characterization, and Bayesian inference methods, and it will include a worked case study based on an analyzed TESS target. This first implementation pass establishes the MNRAS-formatted structure, section layout, and bibliography plumbing required for a full scientific draft.
\end{abstract}

\begin{keywords}
exoplanets -- methods: data analysis -- techniques: photometric -- techniques: radial velocities -- stars: planetary systems
\end{keywords}

\section{Introduction}
The TESS exoplanet pipeline provides a step-by-step workflow for identifying and characterizing transiting exoplanets from TESS light curves.

\section{Methods}
This section will document the target resolution, light-curve ingestion, preprocessing, period search, stellar characterization, and Bayesian transit modeling implemented in the codebase.

\section{Analysis}
This section will connect the software workflow to a representative target, summarize the intermediate products, and describe the diagnostics produced by the pipeline.

\section{Results}
This section will present the period estimates, transit model parameters, stellar properties, and derived planetary properties from the case study.

\section{Discussion}
This section will compare the pipeline outputs with the broader exoplanet literature and discuss limitations, systematics, and future extensions.

\section{Conclusion}
The manuscript will conclude with a summary of the pipeline capabilities and the scientific value of the case-study analysis.

\bibliographystyle{mnras}
\bibliography{references}

\end{document}
